% ---------------------------------------------------------------------------
% Requires: \usepackage{booktabs} \usepackage{graphicx}
% ---------------------------------------------------------------------------

\section{MAEMM expression of SAE Features}
\label{app:verb}

For a held-out SAE feature $f$ we inject $v = \mathrm{unit}(W_{\mathrm{enc}}[:,f])$ at layer~1 on a
marker token ($h \leftarrow h + \lVert h \rVert \, v$), generate best-of-$b$, then re-read each
generated text through the \emph{clean base model with the adapter disabled}. Disabling the adapter
at read time is load-bearing: the model cannot satisfy the metric by encoding something only its own
LoRA recognises. Every generated text is scored against all $65{,}536$ features rather than only its
paired one, which yields each feature's activation on text generated for \emph{other} features at no
extra cost. That control is necessary here: a direction-agnostic baseline --- the same RL fed random
isotropic directions --- reaches approximately the same aggregate score, so the metric is largely
gameable by generic fluent text.

\subsection{Criteria}
\label{app:verb-measure}

Everything is a threshold on one number, and the threshold is not innocuous. The inherited absolute
bar of $1.0$ has no derivation: it spans $1.8\%$ to $561\%$ of a feature's own corpus peak, and
$9.2\%$ of features have a corpus peak below it. $\mathrm{norm\_act}$ is unimodal (mode $0.51$, no
trough), so there is no natural cut either.

\begin{table}[htbp]
\centering
\begin{tabular}{llcc}
\toprule
Definition & What it asks & 27B RL & 8B \\
\midrule
$\mathrm{best\_act} > 1.0$      & inherited absolute bar                        & 0.33 & 0.26 \\
$\mathrm{norm\_act} \geq 0.10$  & $\geq 10\%$ of the feature's corpus peak      & 0.04 & 0.38 \\
$\geq$ top-16 example           & would place in the feature's top-16 examples  & 0.79 & 0.80 \\
$\geq$ weakest top example      & as strong as the weakest known example        & 0.66 & 0.76 \\
$>$ own null p95                & beats $f$'s activation on \emph{other} features' texts & --- & 0.29 \\
\bottomrule
\end{tabular}
\caption{Failure rate under five thresholds on the same underlying number. Base rates move
$20\times$; the rarity relation survives all of them (AUC 0.67--0.91). The 27B arms have no null
control, so those figures are uncontrolled for common features firing on any fluent text.}
\label{tab:verb-criteria}
\end{table}

\begin{figure}[htbp]
\centering
\includegraphics[width=\linewidth]{appendix/eval-inversion/fig6_8b_unverbalized_by_criterion.pdf}
\caption{Failure rate against corpus rarity under five definitions, Qwen3-8B, $n=512$ held-out
features. The purple curve is the only one with a per-feature negative sample. The vertical spread
at fixed rarity is the size of the reporting choice; the common downward slope is the effect.}
\label{fig:verb-criteria-8b}
\end{figure}

\subsection{Rare-feature training}
\label{app:verb-rare}

Strong targets do exist for rare features: mining $10$M tokens gives every one of $9{,}972$ TRAIN
features a full top-16, and the rarest band's rank-0 median activation is $84.3$ against a whole-SAE
median corpus peak of ${\sim}103$. The old bank's problem was scan budget, not corpus scarcity. We
LoRA-finetune on the resulting $159{,}552$ pairs and evaluate on $1{,}520$ features never trained on.

\begin{table}[htbp]
\centering
\begin{tabular}{lccc}
\toprule
Criterion & SFT init & RL & After rare-training \\
\midrule
$>$ own \texttt{null\_p95} (permissive)        & 0.270 & \textbf{0.514} & 0.458 \\
$>$ own \texttt{null\_max} (promiscuity-proof) & 0.100 & \textbf{0.235} & 0.178 \\
$\geq$ top-16 corpus example (strict)          & 0.030 & \textbf{0.074} & 0.046 \\
\bottomrule
\end{tabular}
\caption{Held-out TEST. Rare-feature training improves on its own initialization and generalizes to
features never trained on, but does not reach the RL checkpoint, which received no rare-feature
supervision. It also makes the model more promiscuous: median signal-to-null falls $6.09 \rightarrow
3.33$ against RL's $14.76$, which is why the permissive criterion reports $+18.8$ points of transfer
and the strict one $+7.8$.}
\label{tab:verb-rare}
\end{table}

\subsection{What the failures are}
\label{app:verb-failures}

Of the $966$ features no arm activates, $85$ fire on a token only when a near-constant continuation
follows and $14$ have example sets consisting of one repeated boilerplate string. Both are ill-posed
targets rather than misses --- ${\sim}10\%$ of the set, a floor no training removes. Classifying
\emph{every} evaluated feature by its modal peak token, so the comparison is enrichment rather than
prevalence, the failure rate rises monotonically from content toward structure
(\cref{tab:fail-by-token-class}), identically under all three criteria.

\subsection{The failures are expressible}
\label{app:verb-human}

Calling these features inexpressible would be a claim about the features; the measurement only
supports a claim about this inverter. Given a feature's corpus windows, an LLM writes an activating
passage for it readily (\cref{tab:verb-human}), so a textual preimage exists and is findable --- the
limit is identification from the injected direction, not text production.

\begin{table}[htbp]
\centering
\begin{tabular}{lccc}
\toprule
Set & Features & Sentences & Median \% of corpus peak \\
\midrule
Inverter scores 0.00       & 21 & 30 (29 fire) & 89 \\
\emph{No} arm activates    &  8 & 16 (16 fire) & 91 \\
\bottomrule
\end{tabular}
\caption{Hand-written sentences for features the inverter cannot activate. The second row is the
harder set --- no arm fires them under any conditioning --- and includes two purely positional
numeric features, a digit inside a suite number and inside a company registration number, reached at
$98\%$ and $96\%$.}
\label{tab:verb-human}
\end{table}

To separate lack of exposure from something structural, we cluster the $925$ unactivatable features
by their corpus examples and train on $811$ (prose, news, web UI, disclaimers, metadata), holding out
$113$ of entirely unseen types (recipes, sports, contact blocks) with zero overlap.

\begin{table}[htbp]
\centering
\begin{tabular}{lrcc}
\toprule
Training & Features & TRAIN $\mathrm{norm\_act}$ & TEST $\mathrm{norm\_act}$ \\
\midrule
none                          & --- & 0.004 & 0.008 \\
16 spans each                 & 811 & \textbf{0.344} & 0.035 \\
64 spans each                 & 811 & \textbf{0.387} & 0.031 \\
64 spans, decoder-conditioned & 811 & \textbf{0.352} & 0.049 \\
\midrule
decoder-conditioned, \emph{no} training & --- & --- & \textbf{0.181} \\
\bottomrule
\end{tabular}
\caption{Cluster-disjoint transfer. Trained features go from producing nothing to a third of corpus
peak; held-out features stay at median $\mathrm{best\_act}$ $0.00$ with $78\%$ producing zero, and
their null-beating rate falls \emph{below} baseline ($0.221 \rightarrow 0.195$). Quadrupling the
spans improves the fit and slightly worsens transfer. The last row is the sharpest result: decoder
conditioning with no fine-tuning beats every fine-tuned arm on held-out features, so training on the
$811$ costs the model direction-reading it already had.}
\label{tab:verb-transfer}
\end{table}

The MAEMM is trained only on encoder directions, yet reads the decoder row $W_{\mathrm{dec}}[f]$ --- a
different object, encoder--decoder cosine $0.49$--$0.67$ --- zero-shot. On the $155$ features the
encoder handles it reaches $\mathrm{norm\_act}$ $0.828$ against the encoder's $0.965$; on the $113$
it fails, $0.181$ against $0.008$, with $35\%$ beating $\mathrm{null\_max}$ against the encoder's
$0\%$. Nor is the model merely guessing on what it cannot fire: across $96$ rarity-matched features
at $b=32$, self-consistency is indistinguishable between hard and easy groups ($0.757$ vs.\ $0.754$,
$p = 0.93$) against a cross-feature baseline of $0.516$. It produces one confident cluster of text
pointing elsewhere, which is a worse failure for interpretability than not knowing, because the
output still reads like an explanation.

\begin{figure}[htbp]
\centering
\includegraphics[width=\linewidth]{appendix/eval-inversion/fig9_generalization_failure.pdf}
\caption{\textbf{Left:} 46 hand-written sentences over features the inverter cannot activate, scored
through the same clean-base read: median 90\% of corpus peak against the inverter's 0\%.
\textbf{Right:} training on 811 unactivatable features fits those 811 and transfers nothing to 113
held-out types; decoder conditioning without any training (0.181) beats every fine-tuned arm.}
\label{fig:verb-generalization}
\end{figure}

\subsection{Output diversity: measurement}
\label{app:verb-diversity}

\Cref{fig:verb-diversity} scores texts by the mean over all pairs of a feature's texts of the Jaccard
similarity of their token-trigram sets. Texts are re-tokenized with the Qwen3.6-27B tokenizer
without special tokens, so a rollout's trailing end-of-sequence token does not count, and every
source is tokenized the same way. The cross-feature floor is the same score over $20{,}000$ random
pairs of texts written for different features.

\paragraph{Sources.} All on \texttt{2026-09-21\_sae131k\_2k}, a random sample of $2{,}000$ features
of the 131k SAE. The MAEMM is rl-last16 with $8$ rollouts per feature under vLLM and, as an engine
check, $4$ under HF: both give median $0.043$. The corpus source is each feature's $32$ stored
max-activating windows; the LLM is claude-opus-5, $8$ texts on $159$ features. The rare-feature LoRA
was trained for one epoch on $64{,}432$ spans from the rarest band alone. It is an earlier run than
\cref{app:verb-rare}'s, and none of its $4{,}027$ training features is in this set. We read its HF
rollouts, from which its scores are also computed. The two later arms read the same way: the
rare-weighted arm trained on the \texttt{rw10k} draw, none of whose $7{,}963$ features is in this
set; the cluster arm trained on $195$ of this set's failures, which are dropped. NLA-AV
(\texttt{ceselder/qwen3.6-27b-nla-av}) was run on rows $1500$--$1999$ of the set, $8$ explanations
each at its native $200$-token budget and \texttt{--amp raw}, the paper run's setting ($99\%$ parse
as explanations); those rows are a representative slice ($114$--$133$ features per density
quartile). Its \texttt{<explanation>} tags are stripped before embedding. Its vLLM rollouts are excluded: that path attached
the adapter to the untrained base, so they come from a model nobody trained.

\paragraph{Estimator.} The all-pairs mean is unbiased for the expected pairwise score at any number
of texts $k$. Sources with $4$, $8$ and $32$ texts per feature therefore share a centre, and smaller
$k$ only widens the spread. The distribution over pairs is heavy-tailed, most visibly for the corpus:
$73\%$ of its pairs share no trigram, and a few overlapping windows from one document carry its
mean. We report the share of disjoint pairs beside the mean for that reason.

\paragraph{Embedding measures.} \Cref{fig:verb-diversity-embed,tab:verb-diversity-embed} embed every text, unit-normed,
with bge-small-en-v1.5 (the embedder already used by the evaluation pipeline) and, as a check, with
all-mpnet-base-v2, a different architecture and training set. Every source is read at exactly $k=8$
texts per feature, because unlike the Jaccard mean the Vendi score grows with $k$: the MAEMM's
first $8$ vLLM rollouts, the rare-feature LoRA's $8$ HF rollouts, the LLM's $8$ texts, and for the
corpus the feature's top-8 windows, the same windows the LLM was prompted with. The Vendi score is
$\exp\bigl(-\sum_i \lambda_i \log \lambda_i\bigr)$ over the eigenvalues $\lambda_i$ of $K/k$, where
$K$ is the $k \times k$ cosine kernel. Specificity is the median within-feature mean pairwise cosine
minus the mean cosine of $20{,}000$ random pairs of texts written for different features. Coverage
targets are the feature's corpus windows $9$--$32$, dropping any that share $\geq 0.2$ trigram
Jaccard with one of windows $1$--$8$ (stored windows can overlap within a document); a median of
$19$ remain. Coverage is the mean over targets of the cosine to the nearest of the source's eight
texts, minus the same with a random other feature's texts. The two embedders' cosines live on
different scales (mpnet's cross-feature mean is $0.04$--$0.07$, bge-small's $0.45$--$0.49$), so only
orderings are compared across them.

\begin{table}[htbp]
\centering
\small
\begin{tabular}{lcccccc}
\toprule
& \multicolumn{3}{c}{bge-small} & \multicolumn{3}{c}{all-mpnet-base-v2} \\
\cmidrule(lr){2-4}\cmidrule(lr){5-7}
Source & Vendi & specificity & coverage & Vendi & specificity & coverage \\
\midrule
MAEMM rl-last16          & 2.19 & \textbf{0.31} & 0.108 & 3.10 & \textbf{0.58} & 0.205 \\
corpus top-8             & 2.87 & 0.19 & \textbf{0.143} & 4.45 & 0.34 & \textbf{0.287} \\
LLM (opus-5, $n=159$)    & 3.39 & 0.19 & 0.140 & 5.74 & 0.30 & 0.255 \\
NLA-AV ($n=500$)         & \textbf{1.99} & 0.17 & 0.055 & \textbf{3.00} & 0.37 & 0.105 \\
\midrule
rare-feature LoRA        & 4.17 & 0.05 & 0.075 & 7.26 & 0.08 & 0.145 \\
rare-weighted arm        & 4.24 & 0.04 & 0.070 & 7.37 & 0.07 & 0.142 \\
cluster arm ($n=1{,}805$) & 4.42 & 0.02 & 0.051 & 7.58 & 0.03 & 0.112 \\
\midrule
MAEMM, texts cut to 32 tokens & 2.22 & 0.30 & --- & 3.16 & 0.57 & --- \\
\bottomrule
\end{tabular}
\caption{Diversity in meaning, both embedders; medians over features, $k=8$. Bold marks the extreme
of each column among the sources above the first rule: the MAEMM is the most specific, NLA-AV and
the MAEMM the least spread, and the corpus top-8 covers the feature's other windows best. NLA-AV's
coverage is not comparable, since it writes descriptions of text rather than text like the corpus
windows. Below the rule are the three rare-feature training runs, all read on features they were not
trained on: each is spread only because it is barely specific. On the LLM's $159$ features alone the MAEMM's Vendi is $2.32$
(bge-small) and its raw coverage $0.591$ against the LLM's $0.664$.}
\label{tab:verb-diversity-embed}
\end{table}

\paragraph{Robustness.} The MAEMM's median is flat across the set's four quartiles of pool peak
activation ($0.039$--$0.046$), so the score does not track how strongly a feature fires. The
quartiles in the right panel are formed by rank, because the rare-feature LoRA's scores tie heavily
near zero. Length checks for the Jaccard score cut every text to its first $32$ tokens, and to
tokens $8$--$40$ to remove the MAEMM's shared opening words; for the embeddings, generated texts are
cut to their first $32$ tokens and corpus windows to $32$ tokens around their marked peak, since only
$56\%$ of peaks fall in a window's last $32$ tokens and a blind cut would drop them.

\begin{figure}[htbp]
\centering
\includegraphics[width=\linewidth]{appendix/eval-inversion/fig10_diversity_by_model.pdf}
\caption{Diversity in wording. \textbf{Left:} per-feature mean pairwise token-3-gram Jaccard: box is
the IQR, whiskers the 10th--90th percentile, tick and label the median; $k$ is texts per feature.
The diamond is the same source's cross-feature floor (rl-last16's, $5 \times 10^{-5}$, is drawn at the
axis floor). The all-pairs mean is unbiased in $k$, so rows with $k=4$ and $k=32$ share a centre and
differ only in spread. Scores are not length-matched; the length checks are described above. \textbf{Right:} median
best-of-8 $\texttt{norm\_act}$ by quartile of each MAEMM's own Jaccard, $500$ features per point.}
\label{fig:verb-diversity}
\end{figure}
